\documentclass[11pt]{article}
\usepackage[margin=1in]{geometry}
\usepackage{amsmath,amssymb,bm}
\usepackage{booktabs}
\usepackage{hyperref}
\usepackage{microtype}

\newcommand{\ii}{\mathrm{i}}
\newcommand{\Tr}{\operatorname{Tr}}
\newcommand{\Omrf}{\Omega_{\mathrm{rf}}}

\title{Weak-RF Susceptibility Solver\\for the Magnetometry Page}
\author{Alkali Pumping application}
\date{}

\begin{document}
\maketitle

\section{Purpose and phase convention}

The solver calculates the linear response of a selected observable $O$ in the
upper ground hyperfine manifold to a weak, linearly polarized RF field.  The
RF Hamiltonian is defined as
\begin{equation}
  \frac{H_{\mathrm{rf}}(t)}{\hbar}
  =\Omrf F_i\cos(\omega t),
  \label{eq:rf-hamiltonian}
\end{equation}
where $i\in\{x,y,z\}$ is the selected laboratory RF axis,
$\omega=2\pi f_{\mathrm{rf}}$, and $\Omrf$ is an angular Rabi frequency.
The density matrix is expanded directly about zero RF drive:
\begin{equation}
  \rho(t)=\rho_0+\Omrf\rho^{(1)}(t)+\mathcal O(\Omrf^2).
\end{equation}
Thus the calculation does not assign an arbitrary finite value to
$\Omrf$.

The real cosine drive contains two Fourier components,
\begin{equation}
  \cos(\omega t)=\frac12e^{-\ii\omega t}+\frac12e^{+\ii\omega t}.
\end{equation}
We use $\chi_+(\omega)$ for the coefficient of $e^{-\ii\omega t}$ in
$\partial\langle O\rangle/\partial\Omrf$.

\section{Basis, drive, and readout}

Let $|a\rangle=|F,m\rangle$ denote a ground-state basis vector.  The plotted
RF observable belongs to the upper ground manifold
\begin{equation}
  F=F_+\equiv I+\frac12.
\end{equation}
The available observables are spin components $F_j$ and the off-diagonal
alignment operator
\begin{equation}
  Q_{ij}=\frac12(F_iF_j+F_jF_i).
\end{equation}
For the $Q_{ij}$ option, the indices are chosen cyclically from the RF axis:
\begin{equation}
  x\rightarrow yz,\qquad y\rightarrow zx,\qquad z\rightarrow xy.
\end{equation}
The laboratory operators are expressed in the local frame whose $z$ axis is
the selected quantization axis.

For any readout operator $O$,
\begin{equation}
  \langle O\rangle=\Tr(\rho O)=\sum_{a,b}\rho_{ab}O_{ba}.
  \label{eq:trace}
\end{equation}

\section{Independent adjacent-transition response}

Consider the adjacent pair
\begin{equation}
  |a\rangle=|F,m\rangle,\qquad |b\rangle=|F,m-1\rangle.
\end{equation}
Its steady-state population difference is
\begin{equation}
  D_m=p_a-p_b.
\end{equation}
The signed transition angular frequency is
\begin{equation}
  \omega_m=2\pi\nu_{ab},
\end{equation}
where $\nu_{ab}$ includes the static-field and diagonal light-shift
contributions.  The adjacent-coherence decay rate is
\begin{equation}
  \Gamma_m=\Gamma_{\mathrm{OP},m}
           +\Gamma_{\mathrm{ER},m}
           +\Gamma_{\mathrm{SE},m}.
  \label{eq:gamma}
\end{equation}

Define the RF drive matrix elements
\begin{equation}
  d_{ab}=\langle a|F_i|b\rangle,\qquad
  d_{ba}=\langle b|F_i|a\rangle.
\end{equation}
The coefficients of $e^{-\ii\omega t}$ in the two ordered coherences, per
unit $\Omrf$, are
\begin{align}
  \frac{\partial\rho_{ab,+}}{\partial\Omrf}
  &=\frac{\ii d_{ab}D_m/2}
          {\Gamma_m+\ii(\omega_m-\omega)},
  \label{eq:rhoab}\\
  \frac{\partial\rho_{ba,+}}{\partial\Omrf}
  &=\frac{-\ii d_{ba}D_m/2}
          {\Gamma_m+\ii(-\omega_m-\omega)}.
  \label{eq:rhoba}
\end{align}
The first denominator is the co-rotating resonance and the second is the
counter-rotating contribution.  Both are retained because the applied field
is a real cosine.  In particular, the second term correctly represents a
negative signed transition frequency that resonates at positive drive
frequency.

Using Eq.~\eqref{eq:trace}, the total complex susceptibility is
\begin{equation}
  \chi_+(\omega)=
  \sum_m\left[
    O_{ba}\frac{\partial\rho_{ab,+}}{\partial\Omrf}
   +O_{ab}\frac{\partial\rho_{ba,+}}{\partial\Omrf}
  \right].
  \label{eq:transition-sum}
\end{equation}
Only adjacent coherences within $F_+$ that have nonzero drive and readout
matrix elements contribute.

\section{Full coherence-space solver}

The coupled solver retains every ordered off-diagonal coherence.  Write
\begin{equation}
  \bm c=(\rho_{a_1b_1},\rho_{a_2b_2},\ldots)^T,
  \qquad a_k\ne b_k.
\end{equation}
Its linearized equation of motion is
\begin{equation}
  \frac{d\bm c}{dt}=G\bm c
  +\Omrf\left(\bm s e^{-\ii\omega t}
  +\bm s_-e^{+\ii\omega t}\right).
  \label{eq:eom}
\end{equation}
For the $e^{-\ii\omega t}$ component, set
\begin{equation}
  \bm c(t)=\Omrf\bm c_+(\omega)e^{-\ii\omega t}+\cdots.
\end{equation}
Substitution into Eq.~\eqref{eq:eom} gives
\begin{equation}
  (-\ii\omega I-G)\bm c_+=\bm s,
\end{equation}
and therefore
\begin{equation}
  \boxed{\bm c_+(\omega)=(-\ii\omega I-G)^{-1}\bm s.}
  \label{eq:coherence-solution}
\end{equation}

For an ordered coherence $(a,b)$, the RF source and observable readout are
\begin{align}
  s_{ab}&=\frac{\ii}{2}(F_i)_{ab}(p_a-p_b),
  \label{eq:source}\\
  r_{ab}&=O_{ba}.
  \label{eq:readout}
\end{align}
Combining Eqs.~\eqref{eq:coherence-solution} and \eqref{eq:readout} yields
the implemented resolvent expression
\begin{equation}
  \boxed{
  \chi_+(\omega)=\bm r^T(-\ii\omega I-G)^{-1}\bm s.}
  \label{eq:resolvent}
\end{equation}
Here $T$ means the algebraic row-vector contraction used by the density-matrix
readout.  The necessary complex conjugate relationships are already present
in the Hermitian operator matrix elements and in the ordered-coherence basis.

\subsection{Local generator}

For one alkali species, let
\begin{equation}
  \nu_a=\nu_{\mathrm{LS},a}+\nu_{B,a}
\end{equation}
be the light-shift plus magnetic frequency assigned to state $a$.  The bare
coherence generator for $(a,b)$ contains
\begin{equation}
  G^{(0)}_{ab,ab}
  =-\frac{\Gamma_{\mathrm{OP},a}+\Gamma_{\mathrm{OP},b}}{2}
   -\ii,2\pi(\nu_a-\nu_b).
  \label{eq:bare-generator}
\end{equation}
Electron randomization and spin exchange are represented by linear coherence
maps $M_{\mathrm{ER}}$, $M_{\mathrm{SE,self}}$, and
$M_{\mathrm{SE,cross}}$.  The local generator is
\begin{align}
  G_{\mathrm{loc}}={}&G^{(0)}
  +R_{\mathrm{ER}}(M_{\mathrm{ER}}-I)
  +R_{\mathrm{SE,self}}(M_{\mathrm{SE,self}}-I)\nonumber\\
  &+R_{\mathrm{SE,cross}}(M_{\mathrm{SE,cross}}-I).
  \label{eq:local-generator}
\end{align}
These maps are the derivatives of the corresponding density-matrix collision
maps about the calculated steady state.  Thus their off-diagonal entries can
transfer perturbations between different coherences rather than merely adding
a scalar linewidth.

\subsection{Two-species block generator}

For active alkalis A and B, concatenate their ordered coherences:
\begin{equation}
  \bm c=\begin{pmatrix}\bm c_A\\\bm c_B\end{pmatrix}.
\end{equation}
The joint generator is
\begin{equation}
  G=\begin{pmatrix}
  G_A & R_{\mathrm{SE,cross},A}M_{A\leftarrow B}\\
  R_{\mathrm{SE,cross},B}M_{B\leftarrow A} & G_B
  \end{pmatrix}.
  \label{eq:block-generator}
\end{equation}
For the Alkali-A curve, only A is directly driven and read out:
\begin{equation}
  \bm s^{(A)}=\begin{pmatrix}\bm s_A\\0\end{pmatrix},
  \qquad
  \bm r^{(A)}=\begin{pmatrix}\bm r_A\\0\end{pmatrix}.
\end{equation}
The B coherences can nevertheless modify the A response through the two
off-diagonal spin-exchange blocks.  The B curve is defined conversely.

\section{In-phase, quadrature, and amplitude}

The real time-domain susceptibility is defined with the RF cosine as its phase
reference:
\begin{equation}
  \frac{\partial\langle O(t)\rangle}{\partial\Omrf}
  =X(\omega)\cos(\omega t)+Y(\omega)\sin(\omega t).
  \label{eq:xy-definition}
\end{equation}
Since $\chi_+$ is the coefficient of $e^{-\ii\omega t}$,
\begin{align}
  X(\omega)&=2\operatorname{Re}\chi_+(\omega),
  \label{eq:x}\\
  Y(\omega)&=2\operatorname{Im}\chi_+(\omega),
  \label{eq:y}\\
  A(\omega)&=\sqrt{X(\omega)^2+Y(\omega)^2}
             =2|\chi_+(\omega)|.
  \label{eq:a}
\end{align}
Thus $X$ is the in-phase component, $Y$ is the quadrature component, and $A$
is the nonnegative amplitude.  No sign reversal is applied internally.

The optional \emph{Add $\pi$} controls are display transformations only:
\begin{equation}
  X_{\mathrm{plot}}=-X
  \quad\hbox{or}\quad
  Y_{\mathrm{plot}}=-Y.
\end{equation}
They do not change $\chi_+$, the generator, or the amplitude.

\section{Units and optional display factors}

The spin and alignment operators are represented in dimensionless angular
momentum units.  Because the response is differentiated with respect to the
angular Rabi frequency $\Omrf$ in radians per second, the unscaled
susceptibility has the corresponding inverse-frequency unit (seconds for a
dimensionless spin observable).

Two optional display factors may be applied after solving:
\begin{equation}
  (A,X,Y)_{\mathrm{plot}}
  =C(A,X,Y),
  \qquad
  C=\begin{cases}
  1,&\text{no optional factor},\\
  \Gamma_{\mathrm{ref}},&\text{relaxation normalized},\\
  n,&\text{density factored},\\
  \Gamma_{\mathrm{ref}}n,&\text{both selected}.
  \end{cases}
\end{equation}
Here $\Gamma_{\mathrm{ref}}$ is the adjacent-coherence relaxation rate selected
by the application's reference rule and $n$ is the alkali number density.
These factors do not alter the underlying linear-response solution.

\section{Computational form}

The implementation diagonalizes $G$ once,
\begin{equation}
  G=V\Lambda V^{-1},
\end{equation}
and evaluates the response at every requested frequency as
\begin{equation}
  \chi_+(\omega)
  =\sum_k
  \frac{(\bm r^TV)_k(V^{-1}\bm s)_k}
       {-\ii\omega-\lambda_k}.
  \label{eq:modal}
\end{equation}
A typical eigenvalue is $\lambda_k=-\Gamma_k-\ii\omega_k$.  Equation
\eqref{eq:modal} makes the physical structure clear: each coherence mode
contributes a complex Lorentzian centered near its precession frequency and
weighted by how strongly the RF drive excites it and the selected observable
detects it.

\section{Summary of the solver}

\begin{enumerate}
  \item Calculate the zero-RF steady-state populations and diagonal state
        frequencies.
  \item Construct all ordered coherences and the generator $G$, including
        optical relaxation, electron randomization, and spin exchange.
  \item Construct the source vector $\bm s$ from RF matrix elements and
        population differences.
  \item Construct the readout vector $\bm r$ from the chosen $F_j$ or
        $Q_{ij}$ observable.
  \item Evaluate $\chi_+(\omega)=\bm r^T(-\ii\omega I-G)^{-1}\bm s$.
  \item Convert the phasor to $X=2\Re\chi_+$,
        $Y=2\Im\chi_+$, and $A=2|\chi_+|$.
  \item Apply only the user-selected display factors and optional $\pi$ phase
        flips.
\end{enumerate}

\end{document}
